\section{Introduction}
\label{sec:intro}

Training-free referring image segmentation couples a frozen segmenter that
proposes masks with a frozen vision-language model that scores them, and
returns the best-scoring mask. The field has spent its effort on the scorer,
and the practice of reporting only end metrics has hidden which component
actually fails. We measure it.

Our first finding is negative and decisive: at this operating point the
proposal pool is \emph{not} the bottleneck. On full RefCOCO val the pool
already contains a mask with $\mathrm{IoU}\ge0.5$ for \oracleHitPct\,\% of
expressions, and an oracle that always picked the best member would reach
\oracleRefcoco\,mIoU; the released pipeline returns
\refcocoReleased\,mIoU. The recoverable IoU is thus overwhelmingly in
\emph{selection}, and our error taxonomy attributes it to wrong-instance
selection rather than to mask granularity: on wrong-instance failures the
returned mask is nearly disjoint from the referent (mean IoU around $0.01$)
while a correct mask sat in the pool.

Our second finding is that this failure is a \emph{localization} failure, not
a representation failure. Both quantities the released pipeline uses to rank
candidates are region--text similarities: they compare a whole-region
embedding against a whole-sentence embedding. On expressions that are
descriptive rather than spatial --- RefCOCO+ is constructed without location
words, RefCOCOg uses long relational descriptions --- such scores are
ambiguous across instances of the same category, and the pipeline returns the
wrong one. Replacing the encoder with a stronger one (SigLIP-so400m) makes
selection \emph{worse}, which rules out embedding quality as the deficit.

What the pipeline lacks is an explicit answer to ``where is the referent''.
We obtain one from a frozen vision-language model, asked in natural language
for the referent's bounding box, and use it to select among the candidates the
released scores already propose: the candidate whose mask best fills the box.
The choice of prompt is the whole effect. Asking with the parsed noun phrase
recovers little; asking with the \emph{full expression} --- so that attributes
and relations are part of the query --- is worth
$+13.4$/$+10.6$/$+11.6$\,mIoU on the three benchmarks in a held-out study
(Table~\ref{tab:studies}). The component is local: a cached 7B checkpoint,
$15$\,GB, no API key, no network at inference, and it runs 8-bit quantised at
$9.4$\,GB when GPUs are scarce.

We report the result under the protocol we would want applied to ourselves.
All hyperparameters --- five scalars --- are chosen by Optuna
(\studyTrials{} trials per study) on a held-out split built from RefCOCO
family \emph{train} references, with the objective set to the minimum over
datasets so that no benchmark can be sacrificed; the official validation and
test splits are never used for tuning. Every number in every table is a
re-analysis of full-validation records written by this run, and the tables
themselves are generated from those records by script.

Finally, we report a benchmark where the comparison only makes sense within
class, and we state the boundary to the supervised class exactly. On gRefCOCO,
whose metrics penalise false-positive masks on no-target expressions, the
no-target decision is solvable training-free: a three-prompt presence
ensemble abstains at $N_{\mathrm{acc}}=\wallCombN$ --- above the
\wallTrainedNAcc{} the trained state of the art reports --- and every
same-class (zero-shot) published number is \emph{beaten} on both metrics: the
best published zero-shot cIoU \grefZeroShotCIoU{} (SELF1E-8B) and gIoU
\grefZeroShotGIoU{} (PSALM) fall to \grefCIoU{}/\grefGIoU{}, and our
training-free rule also exceeds the gIoU of the strongest supervised baseline
(InstAlign, \grefPubGIoU). What remains out of reach is the \emph{supervised}
cIoU level (\grefPubCIoU): at the shipped operating point it requires a
targeted realization at $\approx90\%$ of the \emph{best measured oracle}
(pool selection and the construction itself agree at $\approx0.88$), where the
rule delivers $\approx66\%$. By measurement the gap is one of selection
quality --- and within that, one of \emph{which box to ask} rather than which
mask to keep.

Contributions:
\begin{itemize}\itemsep2pt
  \item A measured diagnosis that selection --- specifically instance
  localization --- is the bottleneck of training-free RIS at the released
  operating point, with the pool's ceiling and the error taxonomy quantified
  on full validation splits.
  \item A dense localization prior: a frozen local VLM queried with the full
  expression, used to select among the candidates the released scores
  propose. It is training-free, API-free, and orthogonal to the released
  scoring machinery.
  \item A held-out tuning protocol (\studyTrials{}-trial Optuna studies on
  train-derived refs, minimum-over-datasets objective) and a fully scripted
  evidence chain in which no number is transcribed by hand.
  \item A quantified boundary on gRefCOCO. The no-target decision is
  \emph{solved} above the trained state of the art (N-acc \wallCombN{} vs.\
  \wallTrainedNAcc{}), every same-class (zero-shot) published number is
  \emph{beaten} on both metrics by the training-free rule (\grefCIoU{}/\grefGIoU{}
  vs.\ \grefZeroShotCIoU{}/\grefZeroShotGIoU{}; the strongest supervised gIoU
  reference, \grefPubGIoU, is also exceeded), and the remaining gap to the
  supervised cIoU level is localized by measurement to
  selection quality: that level needs the selector at \wallRealNeedCIoU\,\% of
  the deployable oracle while the rule realizes \wallRealMeas\,\%, with the
  choice of box carrying most of the loss (per-object GT boxes, same rule:
  $0.802$ mean targeted IoU against $0.568$ delivered).
\end{itemize}
